\begin{enumerate}[itemsep=1em]
	\item 
	
	$55~\% \times 160 = 88$
	
	$160-88 = 72$
	
	Il y a donc 88 filles et 72 garçons.
	
	\medskip
	$88\div2=44$
	
	Il y a donc 44 filles qui utilisent les jeux et 44 qui ne les utilisent pas.
	
	\medskip
	$\dfrac{3}{4}\times72=54$
	
	\medskip
	$\dfrac{1}{4}\times72=18$
	
	\medskip
	Il y a donc 54 garçons qui utilisent les jeux et 18 qui ne les utilisent pas.
	
	Finalement, on a : 
	
	\medskip 
	{\renewcommand{\arraystretch}{1.8}
 	\begin{tabularx}{\linewidth}{|X|*{3}{>{\centering\arraybackslash}X|}}
	 	\hline
	 	 & Filles & Garçons & Total\\
	 	 \hline
	 	Utilisent les\newline jeux de cour & 44 & 54 & 98 \\
	 	\hline
	 	N'utilisent pas\newline les jeux de cour & 44 & 18 & 62 \\
	 	\hline
	 	Total & 88 & 72 & 160 \\
	 	\hline
 	\end{tabularx}}
 	
 	\item Ils sont 98 sur un total de 160 élèves à utiliser les jeux de cour, la proportion est donc : $\dfrac{98}{160}=0,6125=\dfrac{61,25}{100}=61,25~\%$.
 	
 	\medskip
 	Il y a environ 61~\% des élèves qui utilisent les jeux de cour.
 	
 	\item Il y a 44 filles qui utilisent les jeux de cour sur un total de 98 élèves.
 	
 	\medskip
 	$\dfrac{44}{98}\approx0,45$ ; soit environ $\dfrac{45}{100}$.
 	
 	\medskip
 	Parmi les élèves ayant déclaré utiliser les jeux de cour, environ 45~\% sont des filles. 

\end{enumerate}